\section{Colimits}
\label{section-colimits}

\noindent
Some of the material in this section overlaps with the general
discussion on colimits in
Categories, Sections \ref{categories-section-limits} --
\ref{categories-section-posets-limits}.
The notion of a preordered set is defined in
Categories, Definition \ref{categories-definition-directed-set}.
It is a slightly weaker notion than a partially ordered set.

\begin{definition}
\label{definition-directed-system}
Let $(I, \leq)$ be a preordered set.
A {\it system $(M_i, \mu_{ij})$ of $R$-modules over $I$}
consists of a family of $R$-modules $\{M_i\}_{i\in I}$ indexed
by $I$ and a family of $R$-module maps $\{\mu_{ij} : M_i \to M_j\}_{i \leq j}$
such that for all $i \leq j \leq k$
$$
\mu_{ii} = \text{id}_{M_i}\quad
\mu_{ik} = \mu_{jk}\circ \mu_{ij}
$$
We say $(M_i, \mu_{ij})$ is a {\it directed system} if $I$ is a directed set.
\end{definition}

\noindent
This is the same as the notion defined in Categories,
Definition \ref{categories-definition-system-over-poset}
and Section \ref{categories-section-posets-limits}.
We refer to Categories, Definition \ref{categories-definition-colimit}
for the definition of a colimit of a diagram/system in any
category.

\begin{lemma}
\label{lemma-colimit}
Let $(M_i, \mu_{ij})$ be a system of $R$-modules over the preordered set $I$.
The colimit of the system $(M_i, \mu_{ij})$ is the quotient $R$-module
$(\bigoplus_{i\in I} M_i) /Q$ where $Q$ is the
$R$-submodule generated by all elements
$$
\iota_i(x_i) - \iota_j(\mu_{ij}(x_i))
$$
where $\iota_i : M_i \to \bigoplus_{i\in I} M_i$
is the natural inclusion. We denote the colimit
$M = \colim_i M_i$. We denote
$\pi : \bigoplus_{i\in I} M_i \to M$ the
projection map and
$\phi_i = \pi \circ \iota_i : M_i \to M$.
\end{lemma}

\begin{proof}
This lemma is a special case of
Categories, Lemma \ref{categories-lemma-colimits-coproducts-coequalizers}
but we will also prove it directly in this case.
Namely, note that $\phi_i = \phi_j\circ \mu_{ij}$ in the above
construction. To show the pair $(M, \phi_i)$ is the colimit we have
to show it satisfies the universal property: for any other such pair
$(Y, \psi_i)$ with $\psi_i : M_i \to
Y$, $\psi_i = \psi_j\circ \mu_{ij}$, there is a unique $R$-module
homomorphism $g : M \to Y$ such that the
following diagram commutes:
$$
\xymatrix{
M_i \ar[rr]^{\mu_{ij}} \ar[dr]^{\phi_i} \ar[ddr]_{\psi_i} & &
M_j\ar[dl]_{\phi_j} \ar[ddl]^{\psi_j} \\
& M \ar[d]^{g}\\
& Y
}
$$
And this is clear because we can define $g$ by taking the
map $\psi_i$ on the summand $M_i$ in the direct sum
$\bigoplus M_i$.
\end{proof}

\begin{lemma}
\label{lemma-directed-colimit}
Let $(M_i, \mu_{ij})$ be a system of $R$-modules over the
preordered set $I$. Assume that $I$ is directed.
The colimit of the system $(M_i, \mu_{ij})$ is canonically
isomorphic to the module $M$ defined as follows:
\begin{enumerate}
\item as a set let
$$
M = \left(\coprod\nolimits_{i \in I} M_i\right)/\sim
$$
where for $m \in M_i$ and $m' \in M_{i'}$ we have
$$
m \sim m' \Leftrightarrow
\mu_{ij}(m) = \mu_{i'j}(m')\text{ for some }j \geq i, i'
$$
\item as an abelian group for $m \in M_i$ and $m' \in M_{i'}$
we define the sum of the classes of $m$ and $m'$ in $M$
to be the class of $\mu_{ij}(m) + \mu_{i'j}(m')$ where
$j \in I$ is any index with $i \leq j$ and $i' \leq j$, and
\item as an $R$-module define for $m \in M_i$ and $x \in R$
the product of $x$ and the class of $m$ in $M$ to be the
class of $xm$ in $M$.
\end{enumerate}
The canonical maps $\phi_i : M_i \to M$ are induced by the canonical
maps $M_i \to \coprod_{i \in I} M_i$.
\end{lemma}

\begin{proof}
Omitted. Compare with
Categories, Section \ref{categories-section-directed-colimits}.
\end{proof}

\begin{lemma}
\label{lemma-zero-directed-limit}
Let $(M_i, \mu_{ij})$ be a directed system.
Let $M = \colim M_i$ with $\mu_i : M_i \to M$.
Then, $\mu_i(x_i) = 0$ for $x_i \in M_i$ if and only if
there exists $j \geq i$ such that $\mu_{ij}(x_i) = 0$.
\end{lemma}

\begin{proof}
This is clear from the description of the directed colimit
in Lemma \ref{lemma-directed-colimit}.
\end{proof}

\begin{example}
\label{example-zero-colimit-different}
Consider the partially ordered set $I = \{a, b, c\}$ with
$a < b$ and $a < c$ and no other strict inequalities.
A system $(M_a, M_b, M_c, \mu_{ab}, \mu_{ac})$
over $I$ consists of three $R$-modules $M_a, M_b, M_c$
and two $R$-module homomorphisms $\mu_{ab} : M_a \to M_b$ and
$\mu_{ac} : M_a \to M_c$.
The colimit of the system is just
$$
M := \colim_{i \in I} M_i = \Coker(M_a \to M_b \oplus M_c)
$$
where the map is $\mu_{ab} \oplus -\mu_{ac}$. Thus the kernel of the
canonical map $M_a \to M$ is $\Ker(\mu_{ab}) + \Ker(\mu_{ac})$.
And the kernel of the canonical map $M_b \to M$ is the image of
$\Ker(\mu_{ac})$ under the map $\mu_{ab}$. Hence clearly
the result of Lemma \ref{lemma-zero-directed-limit} is false for
general systems.
\end{example}

\begin{definition}
\label{definition-homomorphism-directed-systems}
Let $(M_i, \mu_{ij})$, $(N_i, \nu_{ij})$ be
systems of $R$-modules over the same preordered set $I$.
A {\it homomorphism of systems} $\Phi$ from $(M_i, \mu_{ij})$ to
$(N_i, \nu_{ij})$ is by definition a family of $R$-module homomorphisms
$\phi_i : M_i \to N_i$
such that $\phi_j \circ \mu_{ij} = \nu_{ij} \circ \phi_i$
for all $i \leq j$.
\end{definition}

\noindent
This is the same notion as a transformation of functors
between the associated diagrams $M : I \to \text{Mod}_R$
and $N : I \to \text{Mod}_R$, in the language of
categories.
The following lemma is a special case of
Categories, Lemma \ref{categories-lemma-functorial-colimit}.

\begin{lemma}
\label{lemma-homomorphism-limit}
Let $(M_i, \mu_{ij})$, $(N_i, \nu_{ij})$ be
systems of $R$-modules over the same preordered set.
A morphism of systems $\Phi = (\phi_i)$ from $(M_i, \mu_{ij})$ to
$(N_i, \nu_{ij})$ induces a unique homomorphism
$$
\colim \phi_i : \colim M_i \longrightarrow \colim N_i
$$
such that
$$
\xymatrix{
M_i \ar[r] \ar[d]_{\phi_i} & \colim M_i \ar[d]^{\colim \phi_i} \\
N_i \ar[r] & \colim N_i
}
$$
commutes for all $i \in I$.
\end{lemma}

\begin{proof}
Write $M = \colim M_i$ and $N = \colim N_i$ and $\phi = \colim \phi_i$
(as yet to be constructed). We will use the explicit description of $M$
and $N$ in Lemma \ref{lemma-colimit} without further mention.
The condition of the lemma is equivalent to the condition that
$$
\xymatrix{
\bigoplus_{i\in I} M_i \ar[r] \ar[d]_{\bigoplus\phi_i} & M \ar[d]^\phi \\
\bigoplus_{i\in I} N_i \ar[r] & N
}
$$
commutes. Hence it is clear that if $\phi$ exists, then it is unique.
To see that $\phi$ exists, it suffices to show that the kernel of the
upper horizontal arrow is mapped by $\bigoplus \phi_i$ to the kernel
of the lower horizontal arrow. To see this, let $j \leq k$ and
$x_j \in M_j$. Then
$$
(\bigoplus \phi_i)(x_j - \mu_{jk}(x_j))
=
\phi_j(x_j) - \phi_k(\mu_{jk}(x_j))
=
\phi_j(x_j) - \nu_{jk}(\phi_j(x_j))
$$
which is in the kernel of the lower horizontal arrow as required.
\end{proof}

\begin{lemma}
\label{lemma-directed-colimit-exact}
\begin{slogan}
Filtered colimits are exact. Directed colimits are exact.
\end{slogan}
Let $I$ be a directed set.
Let $(L_i, \lambda_{ij})$, $(M_i, \mu_{ij})$, and
$(N_i, \nu_{ij})$ be systems of $R$-modules over $I$.
Let $\varphi_i : L_i \to M_i$ and $\psi_i : M_i \to N_i$ be
morphisms of systems over $I$. Assume that for all $i \in I$ the
sequence of $R$-modules
$$
\xymatrix{
L_i \ar[r]^{\varphi_i} &
M_i \ar[r]^{\psi_i} &
N_i
}
$$
is a complex with homology $H_i$.
Then the $R$-modules $H_i$ form a system over $I$,
the sequence of $R$-modules
$$
\xymatrix{
\colim_i L_i \ar[r]^\varphi &
\colim_i M_i \ar[r]^\psi &
\colim_i N_i
}
$$
is a complex as well, and denoting $H$ its homology we have
$$
H = \colim_i H_i.
$$
\end{lemma}

\begin{proof}
It is clear that
$
\xymatrix{
\colim_i L_i \ar[r]^\varphi &
\colim_i M_i \ar[r]^\psi &
\colim_i N_i
}
$
is a complex. For each $i \in I$, there is a canonical
$R$-module morphism $H_i \to H$ (sending each
$[m] \in H_i = \Ker(\psi_i) / \Im(\varphi_i)$ to the
residue class in $H = \Ker(\psi) / \Im(\varphi)$
of the image of $m$ in $\colim_i M_i$). These give rise
to a morphism $\colim_i H_i \to H$. It remains to
show that this morphism is surjective and injective.

\medskip\noindent
We are going to repeatedly use the description of colimits over $I$
as in Lemma \ref{lemma-directed-colimit} without further mention.
Let $h \in H$.
Since $H = \Ker(\psi)/\Im(\varphi)$ we see that
$h$ is the class mod $\Im(\varphi)$ of an element $[m]$
in $\Ker(\psi) \subset \colim_i M_i$. Choose an
$i$ such that $[m]$ comes from an element $m \in M_i$. Choose
a $j \geq i$ such that $\nu_{ij}(\psi_i(m)) = 0$ which is possible
since $[m] \in \Ker(\psi)$. After replacing $i$ by $j$ and
$m$ by $\mu_{ij}(m)$ we see that we may assume $m \in \Ker(\psi_i)$.
This shows that the map $\colim_i H_i \to H$ is surjective.

\medskip\noindent
Suppose that $h_i \in H_i$ has image zero on $H$. Since
$H_i = \Ker(\psi_i)/\Im(\varphi_i)$ we may represent
$h_i$ by an element $m \in \Ker(\psi_i) \subset M_i$.
The assumption on the vanishing of $h_i$ in $H$ means that
the class of $m$ in $\colim_i M_i$ lies in the image
of $\varphi$. Hence there exists a $j \geq i$ and an $l \in L_j$
such that $\varphi_j(l) = \mu_{ij}(m)$. Clearly this shows that
the image of $h_i$ in $H_j$ is zero. This proves the
injectivity of $\colim_i H_i \to H$.
\end{proof}

\begin{example}
\label{example-colimit-not-exact}
Taking colimits is not exact in general.
Consider the partially ordered set $I = \{a, b, c\}$ with
$a < b$ and $a < c$ and no other strict inequalities,
as in Example \ref{example-zero-colimit-different}.
Consider the map of systems
$(0, \mathbf{Z}, \mathbf{Z}, 0, 0) \to
(\mathbf{Z}, \mathbf{Z}, \mathbf{Z}, 1, 1)$.
From the description of the colimit in
Example \ref{example-zero-colimit-different}
we see that the associated map of colimits is not injective,
even though the map of systems is injective on each object.
Hence the result of Lemma \ref{lemma-directed-colimit-exact}
is false for general systems.
\end{example}

\begin{lemma}
\label{lemma-almost-directed-colimit-exact}
Let $\mathcal{I}$ be an index category satisfying the assumptions of
Categories, Lemma \ref{categories-lemma-split-into-directed}.
Then taking colimits of diagrams of abelian groups over $\mathcal{I}$
is exact (i.e., the analogue of
Lemma \ref{lemma-directed-colimit-exact}
holds in this situation).
\end{lemma}

\begin{proof}
By
Categories, Lemma \ref{categories-lemma-split-into-directed}
we may write $\mathcal{I} = \coprod_{j \in J} \mathcal{I}_j$ with each
$\mathcal{I}_j$ a filtered category, and $J$ possibly empty. By
Categories, Lemma \ref{categories-lemma-directed-category-system}
taking colimits over the index categories $\mathcal{I}_j$ is
the same as taking the colimit over some directed set. Hence
Lemma \ref{lemma-directed-colimit-exact}
applies to these colimits. This reduces the problem to showing that
coproducts in the category of $R$-modules over the set $J$ are exact.
In other words, exact sequences
$L_j \to M_j \to N_j$ of $R$ modules we have to show that
$$
\bigoplus\nolimits_{j \in J} L_j
\longrightarrow
\bigoplus\nolimits_{j \in J} M_j
\longrightarrow
\bigoplus\nolimits_{j \in J} N_j
$$
is exact. This can be verified by hand, and holds even if $J$ is empty.
\end{proof}
